\section{Limitations}

The model scope is limited to the selected compact implementations and one small Chronos variant. The linear temporal proxy lacks the decomposition in published DLinear, and the modern style implementations are not official reproductions. Larger models and broader tuning may produce different results. The current experiment therefore cannot establish a ranking of official architectures or a general limit on neural forecasting.

The protocols also limit comparisons between model groups and feature sets. Each feature set retains its own usable-data split, while Chronos uses a longer univariate context and sampled test origins. The full-origin Chronos check contains only nine unique tickers selected by list order. Its results do not prove that sampled-origin estimates are unbiased across the full universe. Repeated assets and related configurations also prevent us from treating the result rows as independent evidence about a wider population.

The data represent a retrospective snapshot rather than a point-in-time market database. Adjusted-price inputs use the downloaded adjustment factors, whose availability at historical forecast origins remains unverified. The configured asset lists do not reconstruct historical constituents or include every delisted asset, so selection and survivorship effects may matter. Overlapping lists also give some tickers more weight in the aggregate summaries. These limitations restrict claims about historical market-wide performance and real-time deployment.

The corrected batch has stronger provenance than the original archive, but it is not an exact recreation of the original weights. The original archive lacks the complete environment, generation state, prediction traces, and checkpoints needed for that claim. Several corrections and runtime differences accompany the rerun, so we cannot attribute every old--new change to baseline alignment alone. Zero-shot use also does not by itself prove that the pretrained model's training data are wholly separate from this study's assets. Finally, our conclusions concern price-level errors and do not establish results for returns, direction, risk, portfolios, or trading profitability.
